\sct{Coordinate Systems}
\ssc{Coordinate system (座標/坐標系統)}
A coordinate system is a system that uses one or more numbers, or coordinates, to uniquely determine the points on a manifold.
\ssc{Cartesian coordinate system (笛卡爾座標系統) or rectangular coordinate system (直角座標系統)}
\sssc{Definition}
The Cartesian or rectangular coordinate system is a coordinate system in a Euclidean vector space $E$ determined by an orthogonal basis of unit vectors such that the inner product of two vectors in $E$ equals to the dot product of them, and that the norm of $E$ is 2-norm.
\sssc{Lattice points (格子點)}
In a Cartesian coordinate system, vectors whose components are all integers are called lattice points.
\sssc{Quadrant (象限)}
In two dimensions,
\RNum{1}(+, +), \RNum{2}(-, +), \RNum{3}(-, -), \RNum{4}(+, -).
\sssc{Right-hand rule (右手定則)}
Right-hand rule is a common used orientation for a three-dimensional space. In a right-handed Cartesian coordinate system, if you point the thumb of your right hand in the positive $x$-axis direction and your index finger in the positive $y$-axis direction, then your middle finger (extended perpendicularly from the palm) will point in the positive $z$-axis direction. If not specified otherwise, right-handed Cartesian coordinate system is usually used.
\sssc{Octant (卦限)}
In three dimensions, \RNum{1}(+, +, +), \RNum{2}(-, +, +), \RNum{3}(-, -, +), \RNum{4}(+, -, +), \RNum{5}(+, +, -), \RNum{6}(-, +, -), \RNum{7}(-, -, -), \RNum{8}(+, -, -).
\ssc{Polar coordinate system (極座標系)}
\sssc{Definition}
A polar coordinate system is a two-dimensional coordinate system that specifies point positions by a point called pole and a polar axis (a reference ray) that starts from the pole. The two polar coordinates $(r,\theta)$ are:
\bit
\item the radial distance $r\in\bbR_{\geq 0}$ defined as the distance of the point from the origin; and
\item the polar angle, azimuth, or azimuthal angle $\theta$ defined as any signed angle of the ray from the pole to the point from the polar axis with counterclockwise being positive and clockwise being negative.
\eit
For $r<0$, $(r,\theta)$ is defined as $(-r,\theta+\pi)$.
\sssc{Conversion between Cartesian coordinate system}
If the pole and polar axis of a polar coordinate system is the origin and the positive direction of the $x$ axis of a Cartesian coordinate system, then
\[x=r\cos\theta,\]
\[y=r\sin\theta;\]
\[r=\sqrt{x^2+y^2},\]
\[\theta=\operatorname{arctan2}(y,x)+2k\pi,\quad k\in\bbZ.\]
\sssc{Polar rectangle}
A polar rectangle is a region in $\bbR^2$ of the form
\[R=\{(r,\theta)\mid a\le r\le b\land\alpha\le\theta\le\beta\}\]
where $0\le a\le b$ and $0\le\beta-\alpha\le 2\pi$.
\ssc{ISO 31-11 cylindrical coordinate system (圓柱座標系)}
\sssc{Definition}
A cylindrical coordinate system is a three-dimensional coordinate system that specifies point positions by a main axis (a chosen directed line) and an auxiliary axis (a reference ray) with the intersection of the two axes called the origin and the plane passing through the origin and perpendicular to the main axis equipped with polar coordinate system with the polar axis being the auxiliary axis called the reference plane. The three cylindrical coordinates $(\rho,\varphi,z)$ are:
\bit
\item the radial distance $\rho\in\bbR_{\geq 0}$ defined as the perpendicular distance of the point from the main axis;
\item the azimuth or azimuthal angle $\varphi\in(-\pi,\pi]$ defined as the polar angle of the projection of the point on the reference plane and undefined on the main axis; and
\item the axial coordinate or height $z\in\bbR_{\geq 0}$ defined as the signed distance of the point from the reference plane where the sign is defined to be positive if the direction is the same as the main axis and negative if the direction is the negation of the direction of the main axis.
\eit
\sssc{Conversion between Cartesian coordinate system}
If the main axis and the auxiliary axis of a cyclindrical coordinate system is the $z$ axis and the positive direction of the $x$ axis of a Cartesian coordinate system, then
\[x=\rho\cos\varphi,\]
\[y=\rho\sin\varphi,\]
\[z=z;\]
\[\rho=\sqrt{x^2+y^2},\]
\[\varphi=\begin{cases}\sgn(y)\arccos\frac{x}{\sqrt{x^2+y^2}},\quad&y\neq 0\lor x>0\\\pi,\quad&y=0\land x<0\end{cases},\]
\[z=z.\]
\ssc{ISO 80000-2, inclination-azimuth, or physics convention spherical coordinate system (球座標系)}
A spherical coordinate system is a three-dimensional coordinate system that specifies point positions by a polar axis (a chosen directed line) and an auxiliary axis (a reference ray) with the intersection of the two axes called the origin and the plane passing through the origin and perpendicular to the polar axis equipped with polar coordinate system with the polar axis being the auxiliary axis called the reference plane. The three spherical coordinates $(r,\theta,\varphi)$ are:
\bit
\item the radial distance or radius $r\in\bbR_{\geq 0}$ defined as the distance of the point from the origin;
\item the polar angle, inclination, or zenith angle $\theta\in[0,\pi]$ defined as the angle between the ray from the origin to the point and the positive half of the polar axis and undefined at the origin; and
\item the azimuth or azimuthal angle $\varphi\in(-\pi,\pi]$ defined as the polar angle of the projection of the point on the reference plane and undefined on the polar axis (of the coordinate system not the reference plane).
\eit
\sssc{Conversion between Cartesian coordinate system}
If the polar axis and the auxiliary axis of a spherical coordinate system is the $z$ axis and the positive direction of the $x$ axis of a Cartesian coordinate system, then
\[x=r\sin\theta\cos\phi,\]
\[y=r\sin\theta\sin\phi,\]
\[z=r\cos\theta;\]
\[r=\sqrt{x^2+y^2+z^2},\]
\[\theta=\arccos\frac{z}{r},\]
\[\varphi=\begin{cases}\sgn(y)\arccos\frac{x}{\sqrt{x^2+y^2}},\quad&y\neq 0\lor x>0\\\pi,\quad&y=0\land x<0\end{cases}.\]
If the polar axis and the auxiliary axis of a spherical coordinate system is the main axis and the auxiliary axis of a cyclindrical coordinate system, then
\[\rho=r\sin\theta,\]
\[\varphi=\varphi,\]
\[z=r\cos\theta;\]
\[r=\sqrt{\rho^2+z^2},\]
\[\theta=\arctan\frac{\rho}{z},\]
\[\varphi=\varphi.\]